\section{Bölüm sonu çözümlü problemler}

\begin{enumerate}
  \item \textbf{Korelasyon testi.} $r=0{,}50$ ve $n=27$ için sıfır korelasyon testinin $t$ değerini bulun.

  \textbf{Çözüm:} $t=r\sqrt{(n-2)/(1-r^2)}=0{,}50\sqrt{25/0{,}75}\approx2{,}89$'dur. Serbestlik derecesi $n-2=25$'tir.

  \item \textbf{Açıklanan değişkenlik.} $r=-0{,}60$ sonucunu yön ve $r^2$ bakımından yorumlayın.

  \textbf{Çözüm:} İlişki negatiftir; bir değişken arttıkça diğeri doğrusal olarak azalma eğilimindedir. $r^2=0{,}36$ olup basit doğrusal modelde örneklem değişkenliğinin yüzde 36'sı ortak doğrusal değişkenliktir. Bu oran nedensel etki değildir.

  \item \textbf{Yöntem seçimi.} Saçılım grafiğinde ilişki sürekli azalan fakat belirgin eğrisel ve iki uç gözleme duyarlı görünmektedir. Pearson mı Spearman mı önceliklidir?

  \textbf{Çözüm:} Araştırma hedefi tekdüze sıralı ilişkiyse Spearman daha uygundur; doğrusal biçim gerektirmez ve uç büyüklüklerin etkisini azaltır. Yine de grafik ve etkili gözlem incelemesi raporlanmalı, Spearman otomatik bir aykırı değer çözümü sayılmamalıdır.
\end{enumerate}